\documentclass{article}

\usepackage{makecell}

\title{Ch. 15 HW \#4}
\author{Principles of Microeconomics}
\date{\today}

\begin{document}

\maketitle

\section*{Warmup}

\begin{enumerate}

\item Does a competitive firm's price equal its marginal cost in the short run, in the long run, or both? Explain.

\item Draw a diagram showing a U-shaped average total cost curve, a U-shaped average variable cost curve, and a U-shaped marginal cost curve. On it, indicate the price at which the firm makes zero profit or loss.

\end{enumerate}

\section*{Problems and Applications}

\begin{enumerate}

\item Many small boats are made of fiberglass and a resin derived from crude oil. Suppose that the price of oil rises.

	\begin{enumerate}
	
	\item Using diagrams, show what happens to the marginal cost and average total cost curves of an individual boat-making firm and to the market supply curve.
	
	\item What happens to the profits of boat makers in the short run? What happens to the number of boat makers in the long run? 
	
	\end{enumerate}
	
\setcounter{enumi}{2}
\item Consider total cost and total revenue given in the following table:

	\begin{center}
	\begin{tabular}{ c | c | c | c | c | c | c | c | c}
	\textbf{Quantity} & 0 & 1 & 2 & 3 & 4 & 5 & 6 & 7 \\
	\hline
	\makecell{\textbf{Total}\\\textbf{Cost}} & \$8 & 9 & 10 & 11 & 13 & 19 & 27 & 37 \\
	\hline
	\makecell{\textbf{Total}\\\textbf{Revenue}} & \$0 & 8 & 16 & 24 & 32 & 40 & 48 & 56
	\end{tabular}
	\end{center}
	
	\begin{enumerate}
	
	\item Calculate profit for each quantity. How much should the firm produce to maximize profit?
	
	\item Calculate marginal revenue and marginal cost for each quantity. Graph them. (Hint: Put the points between whole numbers. For example, the marginal cost between 2 and 3 should be graphed at 2.5.) At what quantity do these curves cross? How does this relate to your answer to part (a)?
	
	\item Can you tell whether this firm is in a competitive industry? If so, can you tell whether the industry is in a long-run equilibrium?
	
	\end{enumerate}
	
\setcounter{enumi}{4}
\item Suppose the book-printing industry is competitive and begins in a long-run equilibrium.

	\begin{enumerate}
	
	\item Draw a diagram showing the average total cost, marginal cost, marginal revenue, and supply curve of the typical firm in the industry.
	
	\item Hi-Tech Printing Company invents a new process that sharply reduces the cost of printing books. What happens to Hi-Tech's profits and to the price of books in the short run when Hi-Tech's patent prevents other firms from using the new technology?
	
	\item What happens in the long run when the patent expires and other firms are free to use the technology?
	
	\end{enumerate}
	
\setcounter{enumi}{7}
\item The market for fertilizer is perfectly competitive. Firms in the market are producing output but are currently incurring economic losses.

	\begin{enumerate}
	
	\item How does the price of fertilizer compare to the average total cost, the average variable cost, and the marginal cost of producing fertilizer?
	
	\item Draw two graphs, side by side, illustrating the present situation for the typical firm and for the market.
	
	\item Assuming there is no change in either demand or the firms' cost curves, explain what will happen in the long run to the price of fertilizer, marginal cost, average total cost, the quantity supplied by each firm, and the total quantity supplied to the market.
	
	\end{enumerate}

\end{enumerate}

\end{document}